%==============================================================================
% PARTIDA 08: SISTEMA DE PUESTA A TIERRA
% Codigo: OE.5.4.1.1
%==============================================================================
\subsection{OE.5.4.1.1 Puesta a Tierra PAT-1}

\subsubsection{Especificaciones Tecnicas}

Sistema de puesta a tierra tipo vertical (PAT-1) para la vivienda. Debe cumplir con CNE-U Seccion 6 (Puesta a tierra), RNE EM.010 Articulo 15, NTP 370.053, NTP 370.304 y NTP-IEC 60364-5-54.

\begin{itemize}[nosep]
  \item \textbf{Tipo de electrodo:} Varilla vertical de copperweld (acero recubierto de cobre).
  \item \textbf{Diametro:} 5/8'' (15.88 mm) segun NTP 370.053.
  \item \textbf{Longitud:} 2.40 m (segun CNE-U Articulo 600.4, longitud minima 2.40 m).
  \item \textbf{Recubrimiento de cobre:} 254 micras minimo (aplicado por proceso electrolitico).
  \item \textbf{Conexion:} Conector split-bolt de bronce (aleacion Cu-Sn-Zn) para union varilla-conductor.
  \item \textbf{Conductor de tierra:} Cobre electrolitico temple blando, 6 mm2 de seccion (CNE-U Tabla 600.1).
  \item \textbf{Resistencia objetivo:} $R_g < 15\ \Omega$ (valor recomendado para instalaciones residenciales peruanas).
  \item \textbf{Caja de registro:} Concreto prefabricado, 0.30 m x 0.30 m x 0.30 m, con tapa removible y orificio para inspeccion.
  \item \textbf{Gel conductor:} Compuesto electroconductivo a base de carburo de silicio, para relleno del pozo alrededor de la varilla.
\end{itemize}

\paragraph{Calculo de la resistencia de puesta a tierra (metodo de Dwight):}
Para una varilla vertical de longitud $L$ = 2.40 m y diametro $d$ = 0.016 m, en suelo de resistividad $\rho$ = 100 $\Omega$m (valor tipico para suelos arcillosos de Puno):

\begin{equation}
R_g = \frac{\rho}{2\pi L} \ln\left(\frac{4L}{d}\right) = \frac{100}{2\pi \times 2.40} \ln\left(\frac{4 \times 2.40}{0.016}\right) = 6.63 \times \ln(600) = 6.63 \times 6.40 = 42.4\ \Omega
\end{equation}

Este valor supera los 15 $\Omega$ objetivo. Se recomienda:
\begin{itemize}
  \item Utilizar gel conductor (bentonita) para reducir la resistividad efectiva del suelo alrededor de la varilla.
  \item Si es necesario, instalar una segunda varilla paralela separada 2.40 m (la longitud de una varilla) para reducir la resistencia total.
  \item El gel conductor puede reducir la resistividad aparente del suelo en un 40-60\%.
\end{itemize}

\subsubsection{Requisitos de Materiales}

\begin{longtable}{L{5.5cm} L{2.0cm} L{7.0cm}}
\toprule
\textbf{Material} & \textbf{Cant.} & \textbf{Especificacion tecnica} \\
\midrule
Varilla copperweld 5/8'' x 2.40 m & 1 pza | Acero recubierto cobre 254 $\mu$m, con extremo biselado \\
Conector split-bolt bronce 6-10 mm2 & 1 pza | Bronce estañado, para union varilla-conductor \\
Conductor cobre temple blando 6 mm2 desnudo & 25 m | Cobre electrolitico 99.9\%, clase 2 \\
Caja de registro concreto 0.30x0.30x0.30 m & 1 pza | Con tapa removible, con orificio de inspeccion \\
Gel conductor (bentonita o carbon vegetal activado) & 25 kg | Para relleno del pozo alrededor de la varilla \\
Tuberia PVC 3/4'' (proteccion conductor) & 5 m | Para proteger conductor de tierra en tramos expuestos \\
Terminal de compression tipo ojo 6 mm2 & 2 pza | Para conexion a barra de tierra del tablero \\
Pegamento PVC & 1/8 gln | Para sellado de tuberia de proteccion \\
\bottomrule
\end{longtable}

\subsubsection{Metodo de Ejecucion}

\begin{enumerate}
  \item \textbf{Ubicacion:} Seleccionar el punto de instalacion segun plano IE-06, a no menos de 1.00 m de la cimentacion de la vivienda, en zona accesible para inspeccion y mantenimiento.
  \item \textbf{Excavacion del pozo:}
  \begin{enumerate}
    \item Excavar pozo de 0.60 m x 0.60 m de seccion y 2.60 m de profundidad.
    \item En suelos rocosos, la profundidad minima puede reducirse a 1.50 m, pero se debe compensar con mayor uso de gel conductor.
    \item El fondo del pozo debe estar libre de piedras y materiales organicos.
  \end{enumerate}
  \item \textbf{Instalacion de la varilla:}
  \begin{enumerate}
    \item Colocar la varilla en el centro del pozo, con el extremo biselado hacia abajo.
    \item Hincar la varilla mediante golpes con combo de 8 lb, utilizando un protector de varilla (dado de hincado) para no danar el recubrimiento de cobre.
    \item La varilla debe penetrar completamente (2.40 m), dejando 0.10 m sobresaliendo del fondo del pozo para conexiones.
  \end{enumerate}
  \item \textbf{Conexion del conductor:}
  \begin{enumerate}
    \item Pelar 5 cm del conductor de cobre de 6 mm2 (si tiene aislamiento).
    \item Insertar el conductor en el conector split-bolt.
    \item Colocar el conector split-bolt sobre la varilla, a 5 cm del extremo superior.
    \item Apretar el conector con llave de torque a 15 Nm.
    \item Aplicar gel conductor en toda la conexion para evitar corrosion.
  \end{enumerate}
  \item \textbf{Relleno del pozo:}
  \begin{enumerate}
    \item Rellenar el pozo con tierra cernida (sin piedras) mezclada con gel conductor (bentonita) en proporcion 3:1.
    \item Compactar en capas de 0.15 m, humedeciendo ligeramente cada capa.
    \item Los ultimos 0.30 m se rellenan con tierra natural compactada.
  \end{enumerate}
  \item \textbf{Instalacion de caja de registro:}
  \begin{enumerate}
    \item Colocar la caja de registro de concreto sobre el pozo, centrada sobre la varilla.
    \item La tapa debe quedar al ras con el nivel del terreno.
    \item Pasar el conductor de tierra a traves del orificio lateral de la caja.
  \end{enumerate}
  \item \textbf{Tendido del conductor hasta TG-01:}
  \begin{enumerate}
    \item Tender el conductor de 6 mm2 desde la caja de registro hasta la barra de tierra del TG-01.
    \item En tramos expuestos (sobre piso, en muros), proteger con tuberia PVC de 3/4''.
    \item Fijar la tuberia de proteccion con abrazaderas cada 1.00 m.
    \item Conectar a la barra de tierra del TG-01 mediante terminal de compression tipo ojo, con torque de 5 Nm.
  \end{enumerate}
\end{enumerate}

\subsubsection{Control de Calidad}

\begin{longtable}{L{4.5cm} L{4.5cm} L{4.5cm}}
\toprule
\textbf{Parametro} & \textbf{Metodo de verificacion} & \textbf{Criterio de aceptacion} \\
\midrule
Resistencia de puesta a tierra ($R_g$) & Telurómetro (metodo 3 varillas auxiliares, metodo de Wenner) & $R_g < 15\ \Omega$ \\
Continuidad & Multimetro digital entre barra de tierra TG-01 y varilla & < 0.5 $\Omega$ \\
Profundidad de varilla & Inspeccion visual y cinta metrica & 2.40 m ± 0.10 m \\
Conexion split-bolt & Inspeccion visual y prueba de traccion (5 kg) & Firme, sin juego, con gel conductor \\
Caja de registro & Inspeccion visual & Nivelada, tapa accesible, identificada \\
Proteccion de conductor & Inspeccion visual de tuberia PVC en tramos expuestos & Correctamente instalada \\
\bottomrule
\end{longtable}

\paragraph{Protocolo de medicion de resistencia de puesta a tierra:}
\begin{enumerate}
  \item Utilizar telurómetro digital con frecuencia de medicion de 128 Hz.
  \item Metodo de las 3 varillas auxiliares (metodo de caida de potencial, NTP 370.304).
  \item Distancia entre electrodo auxiliar (C2) y electrodo de tension (P2): 62\% de la distancia total.
  \item Distancia total: 20 m (separacion entre varilla PAT-1 y varilla auxiliar C1).
  \item Realizar 3 mediciones y promediar.
  \item Registrar en protocolo: fecha, condiciones climaticas, valor de resistencia, metodo utilizado, operador.
\end{enumerate}

\subsubsection{Metodo de Medicion}

Se medira por juego (jgo) de puesta a tierra completamente instalado, incluyendo varilla, conductor de 6 mm2, caja de registro, conector split-bolt, gel conductor, pozo de tierra, proteccion de conductor y medicion de resistencia verificada.

\subsubsection{Unidad de Medida}

jgo (juego).

\subsubsection{Forma de Pago}

Pago por juego completamente instalado, con medicion de resistencia menor a 15 $\Omega$ verificada por supervision mediante telurómetro calibrado. El protocolo de medicion debe ser entregado como parte de la recepcion de la partida.

%==============================================================================
% PARTIDA 09: ARTEFACTOS DE ALUMBRADO
% Codigo: OE.5.5.1
%==============================================================================
\subsection{OE.5.5.1 Luminaria LED Interior de Adosar 12 W}

\subsubsection{Especificaciones Tecnicas}

Luminaria LED de 12 W para iluminacion general interior, tipo adosar a techo. Debe cumplir con NTP-IEC 60598-1 (Luminarias), NTP 370.304 y EM.010.

\begin{itemize}[nosep]
  \item \textbf{Tipo:} Adosar a techo (surface mounted).
  \item \textbf{Potencia nominal:} 12 W ± 1 W.
  \item \textbf{Flujo luminoso:} Minimo 1200 lm (eficiencia: 100 lm/W).
  \item \textbf{Temperatura de color:} 3000 K -- 4000 K (blanco calido a neutro).
  \item \textbf{Indice de Rendimiento Cromatico (CRI):} > 80.
  \item \textbf{Factor de potencia:} > 0.90.
  \item \textbf{Vida util:} 25000 horas (L70 segun LM-80).
  \item \textbf{Grado de proteccion:} IP-20 minimo (interiores).
  \item \textbf{Material:} Cuerpo de aluminio extrusionado, difusor de policarbonato opal.
  \item \textbf{Driver LED:} Incorporado, con proteccion contra sobrevoltaje (2 kV) y cortocircuito.
  \item \textbf{Tension de entrada:} 220 V AC, 50/60 Hz.
  \item \textbf{Normas:} NTP-IEC 60598-1, NTP-IEC 62031 (Modulos LED), UL 1598.
\end{itemize}

\subsubsection{Requisitos de Materiales}

\begin{longtable}{L{5.0cm} L{2.0cm} L{7.5cm}}
\toprule
\textbf{Material} & \textbf{Cant.} & \textbf{Especificacion tecnica} \\
\midrule
Luminaria LED adosar 12 W & 19 pza | 1200 lm, 3000-4000 K, CRI > 80, FP > 0.90 \\
Tornillo de fijacion 1/4'' x 1'' & 38 pza | Para fijacion a caja octogonal (2 por luminaria) \\
Conector rapido tipo "wago" 3 bornes & 19 pza | Para conexion fase/neutro/tierra (1 por luminaria) \\
Cinta aislante vinilica & 1 pza | Para proteccion de conexiones \\
\bottomrule
\end{longtable}

\subsubsection{Metodo de Ejecucion}

\begin{enumerate}
  \item Verificar la existencia y correcta instalacion de la caja octogonal de salida.
  \item Retirar la tapa ciega de la caja octogonal.
  \item Identificar los conductores: fase (rojo), neutro (blanco), tierra (verde/verde-amarillo).
  \item Fijar la luminaria LED a la caja octogonal mediante los tornillos de fijacion en las orejas de la luminaria.
  \item Conectar los conductores:
  \begin{itemize}
    \item Fase (rojo) al borne L del driver LED.
    \item Neutro (blanco) al borne N del driver LED.
    \item Tierra (verde) al borne de tierra de la luminaria (si aplica).
  \end{itemize}
  \item Colocar el difusor de policarbonato.
  \item Energizar el circuito y verificar el funcionamiento.
\end{enumerate}

\subsubsection{Control de Calidad}

\begin{itemize}
  \item Verificar que la luminaria opere sin parpadeos.
  \item Medir tension en bornes de la luminaria (220 V ± 10\%).
  \item Medir corriente de consumo (debe ser < 0.07 A para 12 W a 220 V).
  \item Verificar correcta fijacion mecanica (sin juego).
\end{itemize}

\subsubsection{Metodo de Medicion}

Se medira por pieza (pza) de luminaria LED instalada y probada.

\subsubsection{Unidad de Medida}

pza (pieza).

\subsubsection{Forma de Pago}

Pago por pieza instalada, probada (encendido verificado) y aprobada por supervision.

%==============================================================================
% PARTIDA 10: EQUIPOS ELECTRICOS Y MECANICOS
% Codigo: OE.5.6.1
%==============================================================================
\subsection{OE.5.6.1 Electrobomba Monofasica de Agua 1 HP, 220V, 60Hz}

\subsubsection{Especificaciones Tecnicas}

Electrobomba centrifuga monofasica de 1 HP para sistema de presurizacion de agua de la vivienda.

\begin{itemize}[nosep]
  \item \textbf{Potencia:} 1 HP (745 W).
  \item \textbf{Tension:} 220 V monofasico, 60 Hz.
  \item \textbf{Corriente nominal:} Aprox. 5.0 A a plena carga (verificar en placa).
  \item \textbf{Proteccion termica:} Incorporada (protector termico de sobrecarga).
  \item \textbf{Caudal:} 40-60 L/min segun altura dinamica total.
  \item \textbf{Altura dinamica total maxima:} 25-35 mca.
  \item \textbf{Conexion:} 1'' de succion y 1'' de descarga (rosca NPT).
  \item \textbf{Control:} Interruptor de nivel automatico (control de nivel por boya electrica o presostato).
  \item \textbf{Proteccion electrica:} Interruptor termomagnetico 2P-16A y tomacorriente servicio pesado 20 A.
\end{itemize}

\subsubsection{Requisitos de Materiales}

\begin{longtable}{L{5.0cm} L{2.0cm} L{7.5cm}}
\toprule
\textbf{Material} & \textbf{Cant.} & \textbf{Especificacion tecnica} \\
\midrule
Electrobomba monofasica 1 HP, 220 V & 1 UND | Centrifuga, con protector termico, 40-60 L/min \\
Interruptor de nivel automatico (boya) & 1 pza | Para control automatico de llenado de tanque \\
Valvula de compuerta 1'' & 2 pza | Para succion y descarga \\
Valvula de retencion 1'' & 1 pza | Para evitar retroceso de agua \\
Nipple galvanizado 1'' x 4'' & 2 pza | Conexion a succion y descarga \\
Union universal galvanizada 1'' & 2 pza | Para desmontaje de la bomba \\
Cinta teflon (sellador de roscas) & 2 pza | Para sellado de conexiones hidraulicas \\
Perno de anclaje 1/2'' x 4'' con tarugo & 4 pza | Para fijacion de la bomba a base de concreto \\
\bottomrule
\end{longtable}

\subsubsection{Metodo de Ejecucion}

\begin{enumerate}
  \item Ubicar la electrobomba en el area tecnica designada (tercer piso), sobre base de concreto de 0.40 m x 0.40 m x 0.15 m.
  \item Fijar la bomba a la base mediante pernos de anclaje de 1/2''.
  \item Conectar tuberia de succion (1'') desde el tanque de agua hasta la entrada de la bomba.
  \item Conectar tuberia de descarga (1'') desde la salida de la bomba hasta la red de distribucion.
  \item Instalar valvula de compuerta en la succion y descarga (para mantenimiento).
  \item Instalar valvula de retencion en la tuberia de descarga (despues de la valvula de compuerta).
  \item Conectar el interruptor de nivel automatico segun instrucciones del fabricante.
  \item Conectar la bomba al tomacorriente de servicio pesado 20 A (circuito dedicado C5).
  \item Verificar que el interruptor termomagnetico asociado sea de 16 A.
\end{enumerate}

\subsubsection{Control de Calidad}

\begin{itemize}
  \item Verificar que la bomba arranque y pare correctamente con el control de nivel.
  \item Medir corriente de operacion con pinza amperimetrica (debe ser menor a la corriente nominal de placa).
  \item Medir tension en bornes: 220 V ± 10\%.
  \item Verificar que no existan fugas de agua en las conexiones hidraulicas.
  \item Verificar que la valvula de retencion funcione correctamente.
\end{itemize}

\subsubsection{Metodo de Medicion}

Se medira por unidad (UND) de electrobomba completamente instalada, incluyendo conexiones electricas, hidraulicas, control de nivel y pruebas de funcionamiento.

\subsubsection{Unidad de Medida}

UND (unidad).

\subsubsection{Forma de Pago}

Pago por unidad instalada, probada en operacion (arranque/parada automatica), sin fugas, y aprobada por supervision.

%==============================================================================
% PARTIDA 11: PRUEBAS ELECTRICAS
% Codigo: OE.5.7.1
%==============================================================================
\subsection{OE.5.7.1 Pruebas de Aislamiento, Continuidad y Resistencia de Puesta a Tierra}

\subsubsection{Especificaciones Tecnicas}

Conjunto de pruebas electricas para verificacion final de la instalacion antes de la puesta en servicio. Deben realizarse segun NTP-IEC 60364-6 (Verificacion de instalaciones electricas), NTP 370.304 y CNE-U Articulo 800.1.

\paragraph{Pruebas obligatorias:}

\begin{longtable}{L{5.0cm} L{5.0cm} L{4.5cm}}
\toprule
\textbf{Prueba} & \textbf{Equipo requerido} & \textbf{Criterio de aceptacion} \\
\midrule
Continuidad de conductores de proteccion & Multimetro digital & Resistencia < 1 $\Omega$ \\
Resistencia de aislamiento & Megohmetro 500 V (hasta 50 V) o 1000 V (> 50 V) & > 1 M$\Omega$ (recomendado > 10 M$\Omega$) \\
Polaridad de tomacorrientes & Comprobador de tomacorrientes (tester) & Correcta L/N/PE \\
Resistencia de puesta a tierra & Telurómetro digital & $R_g < 15\ \Omega$ \\
Funcionamiento de diferenciales & RCD tester o boton TEST de cada ID & Disparo < 40 ms a 30 mA \\
Tension de alimentacion & Multimetro digital true RMS & 220 V ± 10\% (198 V -- 242 V) \\
Corriente de circuito & Pinza amperimetrica & < corriente nominal del interruptor \\
Secuencia de fases & Medidor de secuencia de fases & Correcta (solo para trifasico) \\
\bottomrule
\end{longtable}

\subsubsection{Requisitos de Materiales y Equipos}

\begin{longtable}{L{5.0cm} L{2.0cm} L{7.5cm}}
\toprule
\textbf{Equipo} & \textbf{Cant.} & \textbf{Especificacion tecnica} \\
\midrule
Megohmetro digital 500 / 1000 V & 1 pza | Rango 0-1000 M$\Omega$, precision ±3\% \\
Telurómetro digital & 1 pza | Metodo 3 varillas, rango 0-2000 $\Omega$, frecuencia 128 Hz \\
Multimetro digital true RMS & 1 pza | CAT III 600 V, precision 0.5\% \\
Comprobador de tomacorrientes & 1 pza | Indicacion LED de polaridad L/N/PE \\
RCD tester & 1 pza | Para medicion de tiempo de disparo y corriente de fuga \\
Pinza amperimetrica & 1 pza | Rango 0-200 A AC, CAT III 600 V \\
Protocolo de pruebas & 3 juegos | Formato impreso para registro de resultados \\
\bottomrule
\end{longtable}

\subsubsection{Metodo de Ejecucion}

\begin{enumerate}
  \item \textbf{Antes de energizar:}
  \begin{enumerate}
    \item Verificar que todos los circuitos esten completos y correctamente conectados.
    \item Inspeccion visual general de tableros, cajas, conexiones y canalizaciones.
    \item Prueba de continuidad de todos los conductores de proteccion (tierra).
    \item Medicion de resistencia de aislamiento (F-N, F-T, N-T) con megohmetro de 1000 V.
    \item Verificar que no existan cortocircuitos entre fases, fase-neutro y fase-tierra.
  \end{enumerate}
  \item \textbf{Energizacion:}
  \begin{enumerate}
    \item Conectar el suministro electrico.
    \item Medir tension en la entrada del interruptor general (220 V ± 10\%).
    \item Cerrar el interruptor general.
    \item Medir tension en la barra de distribucion de cada tablero.
  \end{enumerate}
  \item \textbf{Pruebas por circuito:}
  \begin{enumerate}
    \item Cerrar cada interruptor termomagnetico uno por uno.
    \item Verificar polaridad en cada tomacorriente con el comprobador.
    \item Probar cada interruptor de luz (S1 y S3).
    \item Verificar funcionamiento de cada luminaria.
    \item Medir tension en cada tomacorriente.
  \end{enumerate}
  \item \textbf{Pruebas de proteccion diferencial:}
  \begin{enumerate}
    \item Presionar boton TEST de cada interruptor diferencial.
    \item Verificar que dispare inmediatamente (debe cortar el suministro en menos de 40 ms).
    \item Presionar RESET para restablecer.
    \item Si se dispone de RCD tester, medir el tiempo de disparo real.
  \end{enumerate}
  \item \textbf{Medicion de puesta a tierra:}
  \begin{enumerate}
    \item Realizar medicion con telurómetro (metodo de 3 varillas).
    \item Registrar el valor de resistencia.
    \item Si $R_g > 15\ \Omega$, evaluar medidas correctivas (segunda varilla, tratamiento del suelo).
  \end{enumerate}
  \item \textbf{Registro:} Completar el protocolo de pruebas con todos los resultados, fechas, y firma del responsable.
\end{enumerate}

\subsubsection{Control de Calidad}

\begin{itemize}
  \item Todos los equipos de medicion deben tener certificado de calibracion vigente (maximo 1 ano de antiguedad).
  \item Los resultados deben registrarse en el protocolo de pruebas y ser firmados por el supervisor.
  \item Cualquier valor fuera de tolerancia debe ser reportado y corregido antes de la puesta en servicio.
  \item Se debe entregar copia del protocolo de pruebas al propietario.
\end{itemize}

\subsubsection{Metodo de Medicion}

Se medira como un global (glb) que comprende la ejecucion de todas las pruebas descritas para la instalacion completa, incluyendo el uso de equipos de medicion, la elaboracion y entrega del protocolo de pruebas.

\subsubsection{Unidad de Medida}

glb (global).

\subsubsection{Forma de Pago}

Pago por global de pruebas ejecutadas, con protocolo de pruebas presentado, completo, firmado y aprobado por supervision.
